% GENERATED FILE. Edit canonical lessons, then run npm run generate:books.
% canonical-source-hash: fnv1a64:c277a5e0aeb187a6
% canonical-lessons: JA-C502-gekijou, JA-C502-engeki, JA-C502-yakusha, JA-C502-joyuu, JA-C502-enjiru

\chapter{At the Theatre}
\label{ch:ja-at-the-theatre}
\hlchaptermodality{502}

\begin{chapteropening}
\textbf{By the end of this chapter:} \emph{I can talk about plays and the people in them.}

Complete the last lesson of chapter 502: I can talk about plays and the people in them.
\end{chapteropening}

\section[\texorpdfstring{\ja{げきじょう} (gekijō)}{げきじょう (gekijō)}]{\ja{げきじょう} (gekijō) — a theatre}

\label{lesson:JA-C502-gekijou}



\begin{quote}
Before the new one: say the Japanese for a wind chime, then the Japanese for to scatter (something).
\end{quote}

\subsection*{You'll want to know: \ja{げきじょう}}
\textbf{\ja{げきじょう}} — \emph{gekijō} — \textquotedblleft{}a theatre\textquotedblright{}.

\begin{quote}
\textbf{\ja{こんばん} \ja{げきじょうで} \ja{しばいを} \ja{みます。}} — \emph{konban gekijō de shibai o mimasu} — \textquotedblleft{}Tonight I am seeing a play at the theatre.\textquotedblright{}
\end{quote}

\begin{cousinweb}
\textbf{\ja{げきじょう}} (\emph{gekijō}) is Sino-Japanese (kango): two Chinese characters read with their on'yomi, readings borrowed from Chinese. They are \emph{geki}, \textquotedblleft{}drama\textquotedblright{} (Middle Chinese \emph{gjaek}) and \emph{jō}, \textquotedblleft{}location\textquotedblright{} (\emph{drjang}) (KANJIDIC2; Middle Chinese in Baxter's transcription).
\end{cousinweb}

\begin{grammarlens}[title={What you've built}]
One more word for describing what you lived and felt.
\end{grammarlens}

\subsection*{Your turn}
\begin{itemize}
\raggedright
  \item \emph{Say it:} \emph{gekijō}
  \item \emph{Say it:} \emph{gekijō}, once more
  \item \emph{Say it:} \emph{konban gekijō de shibai o mimasu} — Tonight I am seeing a play at the theatre.
  \item \emph{Recall:} say the Japanese for a wind chime, then the Japanese for to scatter (something), then say \emph{gekijō} again
\end{itemize}

\begin{tcolorbox}[breakable,colback=teal!4,colframe=teal!35!black,title={Before you move on}]
What does \ja{げきじょう} mean? (\textquotedblleft{}A theatre\textquotedblright{}.) Say it once more.
\end{tcolorbox}
\section[\texorpdfstring{\ja{えんげき} (engeki)}{えんげき (engeki)}]{\ja{えんげき} (engeki) — drama; the theatre (as an art)}

\label{lesson:JA-C502-engeki}



\begin{quote}
Before the new one: say the Japanese for to scatter (something), then the Japanese for a theatre.
\end{quote}

\subsection*{You'll want to know: \ja{えんげき}}
\textbf{\ja{えんげき}} — \emph{engeki} — \textquotedblleft{}drama; the theatre (as an art)\textquotedblright{}.

\begin{quote}
\textbf{\ja{あねは} \ja{えんげきが} \ja{すきです。}} — \emph{ane wa engeki ga suki desu} — \textquotedblleft{}My elder sister loves the theatre.\textquotedblright{}
\end{quote}

\begin{cousinweb}
\textbf{\ja{えんげき}} (\emph{engeki}) is Sino-Japanese (kango): two Chinese characters read with their on'yomi, readings borrowed from Chinese. They are \emph{en}, \textquotedblleft{}performance\textquotedblright{} (Middle Chinese \emph{yenX}) and \emph{geki}, \textquotedblleft{}drama\textquotedblright{} (\emph{gjaek}) (KANJIDIC2; Middle Chinese in Baxter's transcription).
\end{cousinweb}

\begin{grammarlens}[title={What you've built}]
One more word for describing what you lived and felt.
\end{grammarlens}

\subsection*{Your turn}
\begin{itemize}
\raggedright
  \item \emph{Say it:} \emph{engeki}
  \item \emph{Say it:} \emph{engeki}, once more
  \item \emph{Say it:} \emph{ane wa engeki ga suki desu} — My elder sister loves the theatre.
  \item \emph{Recall:} say the Japanese for to scatter (something), then the Japanese for a theatre, then say \emph{engeki} again
\end{itemize}

\begin{tcolorbox}[breakable,colback=teal!4,colframe=teal!35!black,title={Before you move on}]
What does \ja{えんげき} mean? (\textquotedblleft{}Drama; the theatre (as an art)\textquotedblright{}.) Say it once more.
\end{tcolorbox}
\section[\texorpdfstring{\ja{やくしゃ} (yakusha)}{やくしゃ (yakusha)}]{\ja{やくしゃ} (yakusha) — a stage actor; a player}

\label{lesson:JA-C502-yakusha}



\begin{quote}
Before the new one: say the Japanese for a theatre, then the Japanese for drama.
\end{quote}

\subsection*{You'll want to know: \ja{やくしゃ}}
\textbf{\ja{やくしゃ}} — \emph{yakusha} — \textquotedblleft{}a stage actor; a player\textquotedblright{}.

\begin{quote}
\textbf{\ja{あの} \ja{やくしゃは} \ja{こえが} \ja{いいです。}} — \emph{ano yakusha wa koe ga ii desu} — \textquotedblleft{}That actor has a good voice.\textquotedblright{}
\end{quote}

\begin{cousinweb}
\textbf{\ja{やくしゃ}} (\emph{yakusha}) is Sino-Japanese (kango): two Chinese characters read with their on'yomi, readings borrowed from Chinese. They are \emph{yaku}, \textquotedblleft{}duty\textquotedblright{} (Middle Chinese \emph{ywek}) and \emph{sha}, \textquotedblleft{}someone\textquotedblright{} (\emph{tsyaeX}) (KANJIDIC2; Middle Chinese in Baxter's transcription).
\end{cousinweb}

\begin{grammarlens}[title={What you've built}]
One more word for describing what you lived and felt.
\end{grammarlens}

\subsection*{Your turn}
\begin{itemize}
\raggedright
  \item \emph{Say it:} \emph{yakusha}
  \item \emph{Say it:} \emph{yakusha}, once more
  \item \emph{Say it:} \emph{ano yakusha wa koe ga ii desu} — That actor has a good voice.
  \item \emph{Recall:} say the Japanese for a theatre, then the Japanese for drama, then say \emph{yakusha} again
\end{itemize}

\begin{tcolorbox}[breakable,colback=teal!4,colframe=teal!35!black,title={Before you move on}]
What does \ja{やくしゃ} mean? (\textquotedblleft{}A stage actor; a player\textquotedblright{}.) Say it once more.
\end{tcolorbox}
\section[\texorpdfstring{\ja{じょゆう} (joyū)}{じょゆう (joyū)}]{\ja{じょゆう} (joyū) — an actress}

\label{lesson:JA-C502-joyuu}



\begin{quote}
Before the new one: say the Japanese for drama, then the Japanese for a stage actor.
\end{quote}

\subsection*{You'll want to know: \ja{じょゆう}}
\textbf{\ja{じょゆう}} — \emph{joyū} — \textquotedblleft{}an actress\textquotedblright{}.

\begin{quote}
\textbf{\ja{じょゆうは} \ja{えいがで} \ja{ないていました。}} — \emph{joyū wa eiga de naiteimashita} — \textquotedblleft{}The actress was crying in the film.\textquotedblright{}
\end{quote}

\begin{cousinweb}
\textbf{\ja{じょゆう}} (\emph{joyū}) is Sino-Japanese (kango): two Chinese characters read with their on'yomi, readings borrowed from Chinese. They are \emph{jo}, \textquotedblleft{}woman\textquotedblright{} (Middle Chinese \emph{nrjoX}) and \emph{yū}, \textquotedblleft{}actor\textquotedblright{} (\emph{'juw}) (KANJIDIC2; Middle Chinese in Baxter's transcription).
\end{cousinweb}

\begin{grammarlens}[title={What you've built}]
One more word for describing what you lived and felt.
\end{grammarlens}

\subsection*{Your turn}
\begin{itemize}
\raggedright
  \item \emph{Say it:} \emph{joyū}
  \item \emph{Say it:} \emph{joyū}, once more
  \item \emph{Say it:} \emph{joyū wa eiga de naiteimashita} — The actress was crying in the film.
  \item \emph{Recall:} say the Japanese for drama, then the Japanese for a stage actor, then say \emph{joyū} again
\end{itemize}

\begin{tcolorbox}[breakable,colback=teal!4,colframe=teal!35!black,title={Before you move on}]
What does \ja{じょゆう} mean? (\textquotedblleft{}An actress\textquotedblright{}.) Say it once more.
\end{tcolorbox}
\section[\texorpdfstring{\ja{えんじる} (enjiru)}{えんじる (enjiru)}]{\ja{えんじる} (enjiru) — to play (a role); to perform}

\label{lesson:JA-C502-enjiru}



\begin{quote}
Before the new one: say the Japanese for a stage actor, then the Japanese for an actress.
\end{quote}

\subsection*{You'll want to know: \ja{えんじる}}
\textbf{\ja{えんじる}} — \emph{enjiru} — \textquotedblleft{}to play (a role); to perform\textquotedblright{}.

\begin{quote}
\textbf{\ja{かれは} \ja{えいがで} \ja{いしゃを} \ja{えんじます。}} — \emph{kare wa eiga de isha o enjimasu} — \textquotedblleft{}He plays a doctor in the film.\textquotedblright{}
\end{quote}

Polite \textbf{\ja{えんじます}}, \emph{enjimasu}; an ichidan verb. Past (polite): \textbf{\ja{やくしゃが} \ja{いしゃを} \ja{えんじました。}}, \emph{yakusha ga isha o enjimashita}.

\begin{cousinweb}
\textbf{\ja{えんじる}} (\emph{enjiru}) mixes the two layers, native and Sino-Japanese: Sino-Japanese \emph{en}, \textquotedblleft{}performance\textquotedblright{} (Middle Chinese \emph{yenX}), borrowed from Chinese and the native ending \emph{jiru} (JMdict, KANJIDIC2; Middle Chinese in Baxter's transcription).
\end{cousinweb}

\begin{grammarlens}[title={What you've built}]
One more word for describing what you lived and felt.
\end{grammarlens}

\subsection*{Your turn}
\begin{itemize}
\raggedright
  \item \emph{Say it:} \emph{enjiru}
  \item \emph{Say it:} \emph{enjiru}, once more
  \item \emph{Say it:} \emph{kare wa eiga de isha o enjimasu} — He plays a doctor in the film.
  \item \emph{Recall:} say the Japanese for a stage actor, then the Japanese for an actress, then say \emph{enjiru} again
\end{itemize}

\begin{tcolorbox}[breakable,colback=teal!4,colframe=teal!35!black,title={Before you move on}]
What does \ja{えんじる} mean? (\textquotedblleft{}To play (a role); to perform\textquotedblright{}.) Say it once more.
\end{tcolorbox}
